\clearpage\thispagestyle{empty}
\pdfbookmark[0]{Quellen zum ersten Kometen von 1805}{comet1805-1-sources}
\begin{center}\textbf{Quellen zum ersten Kometen von 1805}\end{center}
\noindent Diese KI-Übertragung in heutiges Deutsch folgt dem Auszug aus Franz Xaver von Zachs Bericht in \href{https://gdz.sub.uni-goettingen.de/id/PPN236007467}{Gauß’ \textit{Werken}, Band~VI} (Göttingen, 1874), S.~265--266. Der Bericht erschien als Beitrag~V in der \textit{Monatlichen Correspondenz}, Band~XIII, Januar 1806, \href{https://www.digitale-sammlungen.de/en/view/bsb10538605?page=87}{S.~79--83}. Auf \href{https://www.digitale-sammlungen.de/en/view/bsb10538605?page=90}{S.~82}--\href{https://www.digitale-sammlungen.de/en/view/bsb10538605?page=91}{83} wird Carl Friedrich Gauß’ Brief vom 5. Dezember zitiert. Er behandelt die Beobachtungen von Johann Sigismund Gottfried Huth, Jacques-Joseph Thulis und Wilhelm Olbers sowie die von Friedrich Wilhelm Bessel berechneten Elemente. Die Beobachtungen von Thulis in Marseille stehen im Novemberheft 1805, Band~XII, auf \href{https://www.digitale-sammlungen.de/en/view/bsb10538604?page=518}{S.~502}. Übertragung und Satz wurden mit KI erstellt. Die Zeitschriftendigitalisate stammen von der Bayerischen Staatsbibliothek.
\clearpage
